% Section 3 -- Related work. Budget: 0.65 page.
% Every contrast is a row of paper/claims.md (section D). Citations in references.bib are verified against
% the ACL Anthology XML (metadata and abstract); citations to references_PLACEHOLDER.bib are UNVERIFIED and
% every claim resting on one carries a \placeholder until it is checked by hand.
\section{Related Work}
\label{sec:related}

\paragraph{Tokenization and sub-word tasks.}
Subword embeddings carry information about the characters of their tokens \citep{2022-naacl-main-179,2022-naacl-main-373}, yet models that know how their tokens are spelled fail to manipulate them: CUTE tests spelling, counting and editing \citep{2024-emnlp-main-177}, EXECUTE extends it to more languages and scripts and to sub-character components of Chinese, Japanese and Korean \citep{2025-findings-acl-95} (\placeholder{whether any of its languages is Vietnamese: page-level check pending}), as does SubTokenTest in practical settings \citep{2026-acl-long-915}. Subword models are brittle to typos \citep{2024-findings-emnlp-86}; byte-level models are more robust on spelling-sensitive tasks \citep{2022-tacl-1-17}. Tokenizer quality and cost differ across languages \citep{2021-acl-long-243,2023-findings-acl-350,petrov-etal-2023-language,2023-emnlp-main-614} (\placeholder{verify the NeurIPS entry}) and across varieties of one language \citep{2025-findings-acl-572}, and the effect of a subword's presence in the vocabulary has been estimated causally with a regression-discontinuity design \citep{2025-acl-long-1374} (see also \citealp{toksuite-2025}, \placeholder{unverified}). Closest to our interventions, re-encodings that keep the characters but change their segmentation lower accuracy across many languages \citep{ghosh-jyothi-2026-noncanonical} (\placeholder{relative drops of 10--24\% as the sweep recorded them; page-level check pending; Vietnamese coverage unconfirmed}), inconsistent tokenizations perplex models on Japanese grammar \citep{2025-acl-short-75}, and full Arabic diacritization fragments tokens and degrades performance \citep{2026-findings-eacl-22}, where diacritics add information; our C1 and C2 arms change only its encoding.

\paragraph{Vietnamese NLP and tokenization.}
Vietnamese encoders and generators segment words before subword tokenization \citep{2020-findings-emnlp-92,N18-5012} or train on large Vietnamese corpora \citep{2022-naacl-srw-18,2311.02945}; Vietnamese LLMs have been evaluated broadly \citep{2024-findings-naacl-182,2025-acl-long-563}, and dialect resources exist for text and speech \citep{2023-findings-emnlp-925,2024-emnlp-main-426,2608.10414}. Tone is encoded by self-supervised speech models \citep{2024-naacl-long-239}. Two recent preprints change the tokenizer instead of diagnosing it, one factorizing Vietnamese syllables into onset, rime and tone before pretraining \citep{2609.21362} (\placeholder{bib entry incomplete: authors}), one comparing tokenizers across Southeast Asian languages \citep{2606.15044}. We know of no study comparing Vietnamese's canonically equivalent encodings on deployed models.

\paragraph{Wordplay and phonological tasks for LLMs.}
Puns were a shared task in 2017 \citep{S17-2005}; LLM pun understanding is shallow and breaks under small reformulations \citep{2024-emnlp-main-657,2025-emnlp-main-1419}; Russian headline wordplay is annotated by type \citep{2025-ranlp-1-15}. Phonological skills are tested in English by grapheme-to-phoneme conversion, syllable counting and rhyme \citep{2024-knowllm-1-1} and in Chinese by homophony, rhyme and phonetic similarity, where models recall pronunciations but struggle to use them \citep{2026-acl-long-1041}; subword tokenization weakens the encoding of rhyme and syllabification, in proportion to the misalignment of token and syllable boundaries \citep{2026-acl-long-634}, the metric our boundary alignment adapts to boundaries inside a syllable. Vietnamese, unlike these, writes onset, rime and tone in the string, so failing to manipulate them is a failure of access, not knowledge. The only computational treatment of nói lái we know of is a detection system \citep{Y18-1063}; the closest LLM analogue is a Filipino syllabification benchmark \citep{pacute-2026-filipino} (\placeholder{unverified}).

\paragraph{Benchmark construction and pre-registration.}
Our tasks follow the minimal-pair logic of BLiMP \citep{2020-tacl-1-25} and symmetric contrastive sets on which a context-blind model scores 50\% \citep{P19-1116}. Static benchmarks saturate and leak \citep{2021-naacl-main-324}; a rule engine mints fresh items. Reporting practice asks for budgets and uncertainty \citep{D19-1224,2411.00640}, and pre-registration has been argued for in NLP \citep{2021-naacl-main-51}; we register hypotheses, metrics, exclusions and analysis before any test-split run.

\paragraph{Probing for sub-word structure.}
Probes show what a representation makes linearly available \citep{P18-1198}, need control tasks \citep{D19-1275}, and show presence, not use \citep{2022-cl-1-7}. Models reconstruct character information in intermediate and higher layers \citep{2025-findings-emnlp-719}; we ask the same of tone, with a structural baseline for what token identity alone predicts, and test use by activation patching \citep{meng-etal-2022-locating,heimersheim-nanda-2024-use}.
